\chapter{Giant Cell Arteritis}
\index{Giant Cell Arteritis}
\label{sec:Giant Cell Arteritis}

\section{Overview}
Giant cell arteritis is the most common primary systemic vasculitis, characterized by granulomatous inflammation of large and medium-sized arteries, with a predilection for the cranial branches of the aorta, particularly the temporal artery, exclusively affecting individuals over 50 years of age (peak incidence 70--80 years), with a female predominance of 2--3:1, and being more common in Northern European descendants.
\section{Causes}
This condition happens because of a Th1- and Th17-driven inflammatory response targeting the vessel wall, with CD4+ T-cell infiltration, multinucleated giant cells, intimal abnormal increase in cells, and luminal occlusion, with IL-6 being a key driver.     
\section{Risk Factors}

\begin{itemize}[noitemsep]
  \item Genetic predisposition and family history
  \item Advancing age
  \item Lifestyle factors (diet, exercise, smoking, alcohol)
  \item Environmental exposures
  \item Underlying medical conditions
  \item Medication use
\end{itemize}
\section{Signs and Symptoms}

\subsection{Common symptoms}
\begin{itemize}[noitemsep]
  \item You may experience new-onset severe headache (the most common symptom, often temporal), scalp tenderness, jaw claudication (pathognomonic), visual disturbances (amaurosis fugax, diplopia, anterior ischemic optic neuropathy leading to irreversible blindness in 15--20\% if untreated), polymyalgia rheumatica symptoms in 40--60\%, fever, weight loss, and an elevated ESR, with aortic involvement (aortitis, aortic aneurysm, aortic dissection) occurring in 10--20\%. 
  \end{itemize}

\subsection{Emergency warning signs}
\begin{itemize}[noitemsep]
  \item Severe or worsening symptoms of giant cell arteritis require immediate medical evaluation
\end{itemize}
\section{Progression and Stages}
The condition may progress through different stages of severity over time. Early stages often involve mild or subtle symptoms that may be overlooked. Moderate stages involve more noticeable symptoms affecting daily function. Advanced stages may involve significant impairment and complications.
\section{Complications}
\begin{itemize}[noitemsep]
  \item Disease progression leading to worsening symptoms
  \item Secondary infections or injuries
  \item Side effects of treatment
  \item Reduced quality of life
  \item Emotional and psychological impact
\end{itemize}
\section{Diagnosis}
To make the diagnosis, doctors look for a high index of suspicion: ESR is almost always markedly elevated ($>$50--100 mm/hr), CRP is elevated, and temporal artery biopsy (TAB, at least 3--5 cm in length) shows granulomatous inflammation with multinucleated giant cells. Comprehensive medical history and physical examination Laboratory tests to assess organ function and rule out other conditions Imaging studies to visualize affected structures Specialized testing depending on the affected system Consultation with relevant specialists Monitoring of disease progression over time

\begin{itemize}[noitemsep]
  \item Comprehensive medical history and physical examination
  \item Laboratory tests to assess organ function
  \item Imaging studies to visualize affected structures
  \item Specialized testing depending on the affected system
  \item Consultation with relevant specialists
\end{itemize}
\section{Prevention}
\begin{itemize}[noitemsep]
  \item Maintain a healthy lifestyle with balanced diet and exercise
  \item Avoid known risk factors
  \item Attend regular medical check-ups for early detection
  \item Follow recommended screening guidelines
\end{itemize}
\section{Treatment Options}
Treatment typically involves with high-dose corticosteroids (prednisone 40--60 mg/day or IV methylprednisolone if vision threatened), with tocilizumab (anti-IL-6 receptor) approved for GCA and used as a steroid-sparing agent. Medications to manage symptoms and address underlying mechanisms Lifestyle modifications and self-care measures Physical therapy and rehabilitation Surgical intervention when indicated Regular monitoring and follow-up care Patient education and support services

\begin{itemize}[noitemsep]
  \item Medications to manage symptoms and address underlying mechanisms
  \item Lifestyle modifications and self-care measures
  \item Physical therapy and rehabilitation
  \item Surgical intervention when indicated
  \item Regular monitoring and follow-up care
\end{itemize}
\section{Living with It}
\begin{itemize}[noitemsep]
  \item Follow your treatment plan as prescribed
  \item Maintain regular follow-up with your healthcare provider
  \item Make necessary lifestyle adjustments
  \item Seek support from family, friends, or support groups
  \item Stay informed about your condition
\end{itemize}
\section{Outlook}
The outlook is good with prompt treatment, though delayed treatment risks irreversible blindness. The outlook varies depending on the specific condition, severity at diagnosis, and response to treatment. Early intervention generally leads to better outcomes. Many conditions can be effectively managed with appropriate treatment. Regular follow-up care is important for monitoring and treatment adjustment.
